\documentclass[11pt,a4paper]{article}
\usepackage[margin=2.2cm]{geometry}
\usepackage[T1]{fontenc}
\usepackage[utf8]{inputenc}
\usepackage{lmodern}
\usepackage{booktabs}
\usepackage{longtable}
\usepackage{tabularx}
\usepackage{enumitem}
\usepackage{fancyhdr}
\usepackage{titlesec}
\usepackage{hyperref}
\usepackage{xcolor}
\usepackage{setspace}
\usepackage{microtype}
\usepackage{listings}

\hypersetup{
  colorlinks=true,
  linkcolor=blue!50!black,
  urlcolor=blue!60!black,
  pdftitle={Open Clinical Record --- Technical Documentation},
  pdfauthor={Open Clinical Record Internship Project}
}

\pagestyle{fancy}
\fancyhf{}
\fancyhead[L]{\small Open Clinical Record}
\fancyhead[R]{\small Technical Documentation}
\fancyfoot[C]{\thepage}
\renewcommand{\headrulewidth}{0.4pt}

\titleformat{\section}{\large\bfseries\color{blue!40!black}}{\thesection}{0.8em}{}
\titleformat{\subsection}{\normalsize\bfseries}{\thesubsection}{0.7em}{}
\setlist[itemize]{leftmargin=1.4em,itemsep=0.25em}
\setlist[enumerate]{leftmargin=1.6em,itemsep=0.25em}
\onehalfspacing

\lstset{
  basicstyle=\ttfamily\small,
  breaklines=true,
  frame=single,
  backgroundcolor=\color{gray!8},
  columns=fullflexible
}

\begin{document}

\begin{titlepage}
\centering
\vspace*{1.5cm}
{\LARGE\bfseries Open Clinical Record}\\[0.4cm]
{\Large Technical Documentation}\\[1.0cm]
{\large Architecture, Implementation, and Operations}\\[0.4cm]
{\large Version 1.0 --- Week 8 Finalization}\\[0.4cm]
{\large 27 September 2026}\\[1.5cm]
\begin{tabular}{ll}
\toprule
\textbf{Repository} & \texttt{elan-thinks/open-clinical-record} \\
\textbf{Branch} & \texttt{main} \\
\textbf{Audience} & Developers, technical mentors, evaluators \\
\textbf{Companion} & SRS v2.0, User Manual \\
\bottomrule
\end{tabular}
\vfill
{\small ASP.NET Core 8 $\cdot$ React + Vite $\cdot$ PostgreSQL $\cdot$ EF Core 8 $\cdot$ JWT}
\end{titlepage}

\tableofcontents
\newpage

\section{Purpose and Scope}

This document describes how Open Clinical Record (OCR) is designed and built. It supports maintenance, evaluation, and handover of the internship MVP.

\textbf{In scope:} patient management, longitudinal patient chart (visits), appointment management including check-in, authentication and RBAC, audit, testing, and local deployment.

\textbf{Out of scope:} pharmacy, laboratory, radiology, billing, FHIR exchange, multi-facility enterprise scheduling, and production multi-tenant SaaS concerns beyond what is needed for a demonstrable MVP.

\section{Architecture}

\subsection{Logical architecture}
\begin{verbatim}
  React SPA (TypeScript, Vite)
        |  HTTPS/HTTP JSON + Bearer JWT
  ASP.NET Core 8 Web API
        |  Application services
        |  EF Core 8 + Npgsql
  PostgreSQL
\end{verbatim}

\subsection{Design principles}
\begin{itemize}
  \item \textbf{Thin controllers}, rich application services for workflows and queries.
  \item \textbf{Authorization at the API}: policies and role claims; the UI only shapes experience.
  \item \textbf{Longitudinal clinical model}: Patient $\rightarrow$ ClinicalVisit $\rightarrow$ content. Prior visits are never overwritten.
  \item \textbf{Explicit state machines} for appointment and visit statuses.
  \item \textbf{Migrations as schema truth}: reference SQL is secondary documentation.
\end{itemize}

\subsection{Major components}

\begin{tabularx}{\textwidth}{@{}p{3.8cm}X@{}}
\toprule
\textbf{Component} & \textbf{Responsibility} \\
\midrule
Auth controller / JWT & Login, token issuance, role claims \\
Patient services & Register, search, update, deceased path \\
AppointmentWorkflowService & Status transitions, reschedule, cancel reason, check-in \\
ClinicalChartService & Chart reads/writes, visit documentation, finalize \\
DashboardService & Aggregate operational statistics \\
AuditService & Persist and query audit events \\
React pages + services & Role-aware UI and typed API clients \\
\bottomrule
\end{tabularx}

\section{Technology Stack}

\begin{tabularx}{\textwidth}{@{}p{3.5cm}X@{}}
\toprule
\textbf{Layer} & \textbf{Choice} \\
\midrule
Frontend & React 18, TypeScript, Vite \\
Backend & ASP.NET Core 8 \\
ORM & Entity Framework Core 8 \\
Database & PostgreSQL \\
Authentication & JWT Bearer \\
Tests & xUnit, WebApplicationFactory \\
CI & GitHub Actions (dotnet test; frontend build) \\
\bottomrule
\end{tabularx}

\section{Repository Layout}

\begin{lstlisting}
open-clinical-record/
  src/backend/OpenClinicalRecord.Api/
  src/frontend/open-clinical-record-web/
  tests/backend/OpenClinicalRecord.Api.Tests/
  docs/
  .github/workflows/
\end{lstlisting}

Backend includes Controllers, Services, Data (DbContext + Migrations), Models/Entities, and DTOs. Frontend includes pages, services (API clients), layouts, routes with role guards, and context for session.

\section{Domain Model}

\subsection{Core relationships}
\begin{verbatim}
Patient (1)
  |-- * Allergies, MedicalHistoryItems
  |-- 0..1 DeathRecord (when Deceased)
  |-- * Appointments
  |-- * ClinicalVisits
         |-- 0..1 VitalSigns
         |-- * Diagnoses
         |-- * ClinicalNotes
\end{verbatim}

Optional link: ClinicalVisit.AppointmentId connects check-in to the visit created for that attendance.

\subsection{Patient status}
Active, Inactive, Deceased. Deceased blocks new appointments and is backed by death-record provenance where recorded.

\subsection{Visit status}
\begin{itemize}
  \item \textbf{Draft} --- editable clinical documentation (default for new visits).
  \item \textbf{Final} --- immutable; further care requires a new visit.
  \item \textbf{Cancelled} --- voided visit.
\end{itemize}
Finalizing a visit that is linked to an appointment sets the appointment to \textbf{Completed} so the patient leaves the active queue.

\subsection{Appointment status and transitions}
Canonical statuses: Scheduled, Waiting, CheckedIn, InProgress, Completed, Cancelled, NoShow.

\begin{tabularx}{\textwidth}{@{}p{2.8cm}X@{}}
\toprule
\textbf{From} & \textbf{Allowed to} \\
\midrule
Scheduled & Waiting, CheckedIn, Cancelled, NoShow \\
Waiting & CheckedIn, Cancelled, NoShow, Scheduled \\
CheckedIn & InProgress, Waiting, Completed, Cancelled \\
InProgress & Completed, CheckedIn \\
Cancelled / NoShow & Scheduled (re-open path) \\
Completed & (terminal) \\
\bottomrule
\end{tabularx}

\vspace{0.4em}
\noindent Cancel requires a non-empty reason. Reschedule is allowed for Scheduled, Waiting, Cancelled, NoShow; blocked for CheckedIn, InProgress, Completed. Check-in creates a Draft visit when appropriate. AppointmentEvent rows preserve history.

\section{Security and Authorization}

\subsection{Authentication}
Users authenticate via login; the API issues a JWT including role claims. Subsequent requests send \texttt{Authorization: Bearer <token>}.

\subsection{Authorization matrix (summary)}

\begin{tabularx}{\textwidth}{@{}Xcccc@{}}
\toprule
\textbf{Action} & \textbf{Admin} & \textbf{Doctor} & \textbf{Nurse} & \textbf{Recep.} \\
\midrule
Register / search patients & Y & Y & Y & Y \\
Book appointment & Y & Y & Y & Y \\
Clinical chart write & N & Y & Y & N \\
Mark deceased & Y & Y & N & Y \\
Clear deceased & Y & Y & N & N \\
Audit / users & Y & N & N & N \\
\bottomrule
\end{tabularx}

\vspace{0.4em}
\noindent Policies such as \texttt{ClinicalStaff}, \texttt{StaffCanBook}, \texttt{CanMarkDeceased}, \texttt{CanClearDeceased}, and \texttt{AdminOnly} encode these rules. Full detail lives in \texttt{docs/05-engineering/access-control-report.md}.

\subsection{Secrets}
Connection strings and JWT signing keys must come from configuration and environment variables, not committed secrets.

\section{API Surface (Illustrative)}

\begin{tabularx}{\textwidth}{@{}p{3.2cm}X@{}}
\toprule
\textbf{Area} & \textbf{Examples} \\
\midrule
Auth & \texttt{POST /api/auth/login} \\
Patients & \texttt{GET/POST /api/patients}, deceased endpoints \\
Appointments & \texttt{GET/POST /api/appointments}, \texttt{PATCH .../status}, \texttt{PATCH .../reschedule} \\
Chart & \texttt{GET/POST/PATCH .../patients/\{id\}/chart/...} \\
Dashboard & \texttt{GET /api/dashboard/stats} \\
Audit & \texttt{GET /api/audit} (Admin) \\
Health & \texttt{GET /api/health}, \texttt{/api/health/ready} \\
\bottomrule
\end{tabularx}

\section{Database and Migrations}

\begin{itemize}
  \item Authoritative schema: \texttt{src/backend/OpenClinicalRecord.Api/Migrations/}
  \item Reference SQL: \texttt{docs/05-data/ocr-complete-database.sql} (may lag; do not prefer over migrations)
  \item Apply with \texttt{dotnet ef database update} (or startup migration if configured)
\end{itemize}

Historical clinical rows use restrictive delete behaviour where appropriate so patient deletion cannot casually erase clinical meaning. Appointment events cascade with appointments. Visit content cascades with the visit. Optional appointment links on visits use SET NULL so removing an appointment does not erase the visit record.

\section{Frontend Structure}
\begin{itemize}
  \item \texttt{src/pages} --- route-level screens
  \item \texttt{src/services} --- API clients and error types
  \item \texttt{src/routes} --- route table and role guards
  \item \texttt{src/layouts} --- application shell
  \item \texttt{src/context} --- authentication session
\end{itemize}
API base URL is configured via \texttt{VITE\_API\_BASE\_URL} (default local API port as documented in the project).

\section{Local Development}

\begin{lstlisting}
# Backend
cd src/backend/OpenClinicalRecord.Api
dotnet run --launch-profile http

# Frontend
cd src/frontend/open-clinical-record-web
npm install
npm run dev
\end{lstlisting}

Typical local URLs: frontend \texttt{http://localhost:5173}, API \texttt{http://localhost:5000}. Development CORS allows the Vite origin.

\section{Testing and Continuous Integration}

\subsection{Automated tests}
Backend integration tests cover:
\begin{itemize}
  \item Access matrix (role permissions)
  \item Appointment workflow (transitions, cancel reason, reschedule, check-in Draft visit)
  \item Clinical documentation and finalize (including InMemory-safe patterns)
  \item Longitudinal visits (history not overwritten)
\end{itemize}

\begin{lstlisting}
cd tests/backend/OpenClinicalRecord.Api.Tests
dotnet test
\end{lstlisting}

\subsection{CI}
GitHub Actions runs backend tests and frontend production build (TypeScript check + Vite build) on pushes to \texttt{main}.

\subsection{Manual verification}
End-to-end checklist: \texttt{docs/09-testing/e2e-clinical-flow-checklist.md}.

\section{Operational Notes}

\begin{itemize}
  \item Prefer environment-based configuration for JWT key length and database credentials.
  \item Training seed users exist only for development demonstration.
  \item Health endpoints support basic liveness and readiness checks.
  \item Backup and restore of PostgreSQL should be defined before any production-like deployment.
\end{itemize}

\section{Known Limitations}
\begin{itemize}
  \item Reference SQL dump may lag the latest EF migrations.
  \item Reports are basic aggregates, not a full analytics platform.
  \item No multi-tenant or multi-facility model.
  \item Frontend automated UI tests are limited relative to backend API tests.
\end{itemize}

\section{Related Documents}
\begin{itemize}
  \item SRS v2.0 --- formal requirements and acceptance criteria
  \item User Manual --- role-based operational guide
  \item Clinical domain rules --- status machines and business rules
  \item Access-control report --- live permission matrix and seed accounts
  \item Project structure --- repository map
\end{itemize}

\section*{Document control}
Technical Documentation version 1.0 (27 September 2026). Implementation evidence is the \texttt{main} branch of the project repository.

\end{document}
